\renewcommand{\abstractname}{Sommario}

\begin{abstract}
   Un terremoto di grande magnitudo può indurre attività vulcanica nella regione circostante.
   Sebbene diversi studi abbiano evidenziato una correlazione tra i due eventi da un punto di vista statistico, i meccanismi fisici alla base non sono stati del tutto compresi.
   Il 29 luglio 2025 un terremoto di magnitudo \(\mathrm{M}\,8.8\) ha colpito le coste della Kamchatka in Russia e nei giorni seguenti due vulcani, il Klyuchevskoy e il Krasheninnikov, hanno eruttato.
   
   In questa tesi si è analizzata la variazione del campo di sforzo nella regione della Kamchatka a seguito del sisma, come possibile causa delle eruzioni.
   Gli sforzi sono stati calcolati tramite un modello a faglia finita basato su dislocazioni rettangolari. Quindi, sono stati simulati possibili percorsi di risalita del magma prima e dopo l'evento sismico. I risultati della simulazione mostrano che, nella zona considerata, il terremoto non ha fornito un contributo sufficiente per cambiare le traiettorie di risalita in modo considerevole.
\end{abstract}
